\documentclass[11pt,a4paper]{article}
\usepackage{curso}
\begin{document}
\portada{Lección}{1.7}{Arco histórico}

\noindent\textbf{Unidad 1.}\quad\textbf{Objetivo:} situar las ideas del curso en su historia y asociar a cada figura
la herramienta que usaremos. \textbf{Prerrequisitos:} lecciones 1.1--1.6.

\section{Tres ríos}
El cómputo termodinámico es la confluencia de (i) la termodinámica, nacida de las máquinas de vapor; (ii) la teoría de
la información, nacida de las telecomunicaciones; y (iii) las redes neuronales inspiradas en la física. A ellos se suma,
desde finales del siglo XX, (iv) la termodinámica estocástica.

\section{Línea de tiempo}
\begin{center}\small
\begin{tabularx}{\linewidth}{llX}\toprule
Año & Figura & Idea y herramienta del curso\\\midrule
1824 & S. Carnot & Eficiencia máxima $\eta=1-T_f/T_c$: los límites dependen solo de temperaturas.\\
1865 & R. Clausius & Segunda ley y entropía, $\dd S=\delta Q_{\text{rev}}/T$.\\
1877 & L. Boltzmann & $S=k\ln W$: la entropía cuenta microestados (lección 1.1).\\
1902 & J. W. Gibbs & Ensambles estadísticos, $p=e^{-\beta E}/Z$ (lección 1.2): el «programa» de un TSU.\\
1948 & C. Shannon & Entropía de la información $H=-\sum p\log_2p$ (lección 1.3).\\
1961 & R. Landauer & $Q\ge\kB T\ln2$ por bit borrado \cite{landauer} (lección 1.5).\\
1973--82 & C. Bennett & Cómputo reversible; resolución del demonio de Maxwell \cite{bennett} (lección 1.4).\\
1982 & J. Hopfield & Redes con $E(s)=-\frac12\sum W_{ij}s_is_j$: memoria asociativa \cite{hopfield} (lección 4.6).\\
1985 & Ackley, Hinton, Sejnowski & Máquina de Boltzmann y su regla de aprendizaje \cite{ackley} (lección 4.7).\\
1997 & C. Jarzynski & $\langle e^{-\beta W}\rangle=e^{-\beta\Delta F}$ fuera del equilibrio \cite{jarzynski} (lección 3.3).\\
1999 & G. Crooks & Teorema de fluctuación de trayectorias (lección 3.3).\\
2024 & Hopfield y Hinton & Premio Nobel de Física por las redes neuronales.\\
2019--26 & varios & Informe CCC sobre cómputo termodinámico \cite{conte}; primeros TSU y SPU (Unidad 5).\\\bottomrule
\end{tabularx}
\end{center}

\section{Ecuaciones emblemáticas}
\[\eta_{\text{Carnot}}=1-\frac{T_f}{T_c},\quad S=k\ln W,\quad p=\frac{e^{-\beta E}}{Z},\quad H=-\sum p\log_2p,\quad Q\ge\kB T\ln2,\quad\langle e^{-\beta W}\rangle=e^{-\beta\Delta F}.\]
\begin{ejemplo}
Carnot entre 600 K y 300 K: $\eta=1-300/600=0{,}5$. Jarzynski implica la segunda ley: por Jensen,
$e^{-\beta\langle W\rangle}\le\langle e^{-\beta W}\rangle=e^{-\beta\Delta F}\Rightarrow\langle W\rangle\ge\Delta F$.
\end{ejemplo}

\begin{aplicacion}
Carnot/Clausius: límites y eficiencia (Unidad 3). Boltzmann/Gibbs: programa del TSU (Unidades 4--5).
Shannon/Landauer/Bennett: costo de la información. Hopfield/Hinton: Unidad 4. Jarzynski/Crooks: energías libres y
determinantes por estadística de trabajo (Unidades 3 y 6).
\end{aplicacion}

\begin{error}
\begin{itemize}[nosep]
\item Shannon no «copió» a Gibbs: la equivalencia formal se reconoció después.
\item La resolución del demonio está en el borrado (Bennett), no en la medición.
\item Las máquinas de Boltzmann no eran erróneas: eran caras de muestrear en hardware digital.
\item Jarzynski no contradice la segunda ley: la implica en promedio.
\end{itemize}
\end{error}

\begin{thebibliography}{9}\small
\bibitem[8]{landauer} R. Landauer, \emph{IBM J. Res. Dev.} 5, 183 (1961).
\bibitem[9]{bennett} C. H. Bennett, Notes on Landauer's principle..., 2003.
\bibitem[12]{jarzynski} C. Jarzynski, \emph{Phys. Rev. Lett.} 78, 2690 (1997).
\bibitem[14]{ackley} D. H. Ackley, G. E. Hinton y T. J. Sejnowski, \emph{Cognitive Science} 9, 147 (1985).
\bibitem[15]{hopfield} J. J. Hopfield, \emph{PNAS} 79, 2554 (1982).
\bibitem[6]{conte} T. Conte et al., \emph{Thermodynamic Computing}, 2019.
\end{thebibliography}
\end{document}
